%% toc: P3 32 32.5,32.6,32.7
\chapter{The Stability Spectrum and Theories Separating the Stability Classes}
\chaptermark{Stability Spectrum and Separating Theories}
\label{sec:p03-w03-p03-032-spectrum}
\index{stability!cardinal spectrum}

\paragraph{Purpose.}
A complete theory \(T\) is stable (\cref{def:p03-stability}) when no
formula has the order property, and it is \(\lambda\)-stable
(\cref{def:p03-stability-spectrum}) when over every parameter set of
cardinality at most \(\lambda\) there are at most \(\lambda\) complete
\(1\)-types.  This chapter proves that the first, purely combinatorial,
condition decides the second at every cardinal \(\lambda\) with
\(\lambda^{\abs T}=\lambda\), and that an unstable theory is
\(\lambda\)-stable for no infinite \(\lambda\).  It proves that stability in
every \(\lambda\geq2^{\abs T}\), Shelah's superstability, is equivalent to
finite local character of forking, through Shelah's degree.  It then constructs, with all
properties proved, two theories that separate the stability classes: one that
is superstable but not \(\omega\)-stable, and one that is stable but not
superstable.  The separation of the stable theories from the simple ones by
the random graph is proved in \cref{thm:p03-w03-p03-031-random-graph} and
\cref{def:p03-w03-p03-031-random-graph-example}.

\paragraph{Inputs and conclusions.}
The input is a complete theory \(T\) in a language \(\Lang\) with
\(\abs T=\max(\abs\Lang,\aleph_0)\), read in the ambient model
\(\mathfrak C\) of \cref{def:p03-w03-p03-031-monster}; for each cardinal
\(\lambda\) under discussion the saturation cardinal \(\bar\kappa\) of that
convention is taken larger than \(2^\lambda\), so that every set of
cardinality at most \(\lambda\) is small.  The set \(S_n(A)\) of complete
\(n\)-types over \(A\) does not depend on this choice.  The conclusions are
the formula-by-formula equivalences of
\cref{thm:p03-w03-p03-032-formula}, the spectrum theorem
\cref{thm:p03-w03-p03-032-spectrum}, the superstability theorem
\cref{thm:p03-w03-p03-032-superstable} and the two separating theories of
\cref{thm:p03-w03-p03-032-tp,thm:p03-w03-p03-032-te}.

\paragraph{Source.}
The development of
\cref{sec:p03-w03-p03-032-ranks,sec:p03-w03-p03-032-shelah} follows the order
of Shelah's 1971 paper \cite[\S~2, pp.~296--305]{Shelah1971Stability}: the
properties of a formula listed there as (A) to (F), their equivalence
\cite[Theorem~2.9, p.~303]{Shelah1971Stability}, the passage from formulas to
theories \cite[Lemmas~2.10--2.12, pp.~303--304]{Shelah1971Stability} and
the equivalences for theories \cite[Theorem~2.13,
pp.~304--305]{Shelah1971Stability}.  The step from many types to the order
property uses the set-theoretic theorem of Erd\H{o}s and Makkai
\cite[Theorem]{ErdosMakkai1966}, which the paper invokes
\cite[Lemma~2.1, p.~298]{Shelah1971Stability} and which is proved here in
full.  Where the paper's rank is defined by an induction on ordinals
\cite[Definitions~2.4--2.5, pp.~299--300]{Shelah1971Stability}, the rank
below is defined through the trees of its Lemma~2.7
\cite[Lemma~2.7, pp.~301--302]{Shelah1971Stability}; the finite values agree
by that lemma, and only finite values are used.

\section{Local ranks and the role of parameter sets}
\label{sec:p03-w03-p03-032-ranks}

A single formula \(\varphi(x;y)\), with the variables split into object
variables \(x\) and parameter variables \(y\), already carries the whole
stability question.  The point of this section is that, when the formula is
stable, a complete type over an arbitrary parameter set \(A\) is pinned down
by a \emph{finite} part of itself; the parameter set only supplies the
instances to be decided.

\begin{hbdefinition}{def:p03-w03-p03-032-phi-types}{Instances, \(\varphi\)-types and \(S_\varphi(A)\)}
Let \(\varphi(x;y)\) be a formula with \(x\) a finite tuple of object
variables and \(y\) a finite tuple of parameter variables.  Write
\(\varphi^1\defeq\varphi\) and \(\varphi^0\defeq\neg\varphi\).  For a set
\(A\), a \emph{\(\varphi\)-instance over \(A\)} is a formula
\(\varphi^\varepsilon(x;b)\) with \(\varepsilon\in\{0,1\}\) and \(b\) a tuple
from \(A\) of the length of \(y\).  A \emph{\(\varphi\)-type} is a set of
\(\varphi\)-instances; it is \emph{complete over \(A\)} when it is consistent
(realized in \(\mathfrak C\)) and contains, for each tuple \(b\) from \(A\),
exactly one of \(\varphi(x;b)\) and \(\neg\varphi(x;b)\).  The set of
complete \(\varphi\)-types over \(A\) is \(S_\varphi(A)\).  For a complete
type \(p\) over \(A\), \(p\mathbin{\upharpoonright}\varphi\) is the set of
\(\varphi\)-instances over \(A\) that belong to \(p\).
\index{phi-type@\(\varphi\)-type}
\end{hbdefinition}

These are the paper's \(\varphi\)-\(m\)-types and \(S^m_\varphi(A)\)
\cite[Definition~2.1, p.~297]{Shelah1971Stability}, with \(m\) the length of
\(x\).  The restriction \(p\mapsto p\mathbin{\upharpoonright}\varphi\) maps
\(S_{\abs x}(A)\) onto \(S_\varphi(A)\), because every consistent set of
formulas over \(A\) extends to a complete type over \(A\).

\begin{hbdefinition}{def:p03-w03-p03-032-tree-rank}{\(\varphi\)-trees and the tree rank}
Let \(\varphi(x;y)\) be as above and \(\pi(x)\) a set of formulas with
parameters from a small set.  For \(n<\omega\), a \emph{\(\varphi\)-tree of
height \(n\) above \(\pi\)} is a family \((b_s:s\in 2^{<n})\) of tuples of
\(\mathfrak C\) indexed by the binary sequences of length below \(n\) such
that, for every \(\eta\in 2^n\),
\[
  \pi(x)\cup\{\varphi^{\eta(k)}(x;b_{\eta\mathbin{\upharpoonright}k}):k<n\}
\]
is realized in \(\mathfrak C\).  A tree of height \(0\) is the empty family,
and exists exactly when \(\pi\) is consistent.  The \emph{tree rank}
\(\mathrm R_\varphi(\pi)\) is \(-1\) if \(\pi\) is inconsistent, the greatest
\(n\) for which a \(\varphi\)-tree of height \(n\) above \(\pi\) exists if
there is a greatest one, and \(\infty\) otherwise.  Write
\(\mathrm R_\varphi(x=x)\) for the rank of the empty set of formulas in the
variables \(x\).
\index{tree rank}
\end{hbdefinition}

The tree rank is the paper's \(\operatorname{Rank}_\varphi\) whenever either
is finite \cite[Lemma~2.7, p.~301]{Shelah1971Stability}.  Its value depends
on the formulas of \(\pi\) and on \(\varphi\), not on a parameter set chosen
to contain them.

\begin{hblemma}{thm:p03-w03-p03-032-rank-basics}{proved=annex:p03-w03-p03-032-rank-basics}{Monotonicity, finite character and splitting}
Let \(\varphi(x;y)\) be a formula and \(\pi,\rho\) sets of formulas in
\(x\) with parameters from a small set.
\begin{enumerate}
  \item If \(\pi\subseteq\rho\), then
        \(\mathrm R_\varphi(\rho)\leq\mathrm R_\varphi(\pi)\).
  \item For each \(n<\omega\), a \(\varphi\)-tree of height \(n\) above
        \(\pi\) exists if and only if one exists above every finite subset of
        \(\pi\).  Hence, if \(\mathrm R_\varphi(\pi)<\infty\), some finite
        \(\pi_0\subseteq\pi\) has
        \(\mathrm R_\varphi(\pi_0)=\mathrm R_\varphi(\pi)\).
  \item If \(\mathrm R_\varphi(\pi\cup\{\varphi(x;b)\})\geq n\) and
        \(\mathrm R_\varphi(\pi\cup\{\neg\varphi(x;b)\})\geq n\), then
        \(\mathrm R_\varphi(\pi)\geq n+1\).
\end{enumerate}
\end{hblemma}

\begin{hblemma}{thm:p03-w03-p03-032-local-character}{proved=annex:p03-w03-p03-032-local-character}{A finite part decides a \(\varphi\)-type}
Let \(\varphi(x;y)\) be a formula with
\(\mathrm R_\varphi(x=x)<\infty\), let \(A\) be small and
\(p\in S_\varphi(A)\).  Put \(k=\mathrm R_\varphi(p)\) and choose a finite
\(p_0\subseteq p\) with \(\mathrm R_\varphi(p_0)=k\).  Then for every tuple
\(b\) from \(A\) of the length of \(y\),
\[
  \varphi(x;b)\in p
  \quad\Longleftrightarrow\quad
  \mathrm R_\varphi\bigl(p_0\cup\{\varphi(x;b)\}\bigr)\geq k .
\]
Consequently \(p\) is determined by \(p_0\), and
\(\abs{S_\varphi(A)}\leq\abs A+\aleph_0\).
\end{hblemma}

The finite set \(p_0\) is the local character of a stable formula in its
plainest form: no matter how large \(A\) is, the decision on every instance
over \(A\) is read off a fixed finite set of instances, by a rank
computation that does not mention \(A\).  Lemma~2.5 of the paper records the
uniqueness behind this
\cite[Lemma~2.5\,(C), (D), p.~300]{Shelah1971Stability}.

\begin{hblemma}{thm:p03-w03-p03-032-product}{proved=annex:p03-w03-p03-032-product}{Types are determined by their local parts}
For every small set \(A\),
\(\abs{S_1(A)}\leq\prod_{\varphi}\abs{S_\varphi(A)}\), where \(\varphi\)
ranges over the formulas \(\varphi(x;y)\) of \(\Lang\) with \(x\) a single
variable; there are at most \(\abs T\) of them.
\end{hblemma}

\begin{hblemma}{thm:p03-w03-p03-032-one-to-many}{proved=annex:p03-w03-p03-032-one-to-many}{One variable suffices}
Let \(\lambda\) be infinite and suppose \(\abs{S_1(A)}\leq\lambda\) for every
\(A\) with \(\abs A\leq\lambda\).  Then \(\abs{S_m(A)}\leq\lambda\) for every
\(m<\omega\) and every \(A\) with \(\abs A\leq\lambda\).
\end{hblemma}

This is \cite[Lemma~2.10, p.~303]{Shelah1971Stability}; it shows that the
definition of \(\lambda\)-stability through \(1\)-types in
\cref{def:p03-stability-spectrum} agrees with the paper's definition through
types of every finite length \cite[Definition~2.1\,(C),
p.~297]{Shelah1971Stability}.

\section{The stability spectrum in Shelah's presentation}
\label{sec:p03-w03-p03-032-shelah}

\subsection{From trees to many types, and from many types to an order}

\begin{hblemma}{thm:p03-w03-p03-032-tree-count}{proved=annex:p03-w03-p03-032-tree-count}{Trees of every height give many types}
Let \(\varphi(x;y)\) be a formula with \(\mathrm R_\varphi(x=x)=\infty\).
Let \(\lambda\) be an infinite cardinal and \(\mu\) the least cardinal with
\(2^\mu>\lambda\).  Then there is a set \(A\) with \(\abs A\leq\lambda\) and
\(\abs{S_\varphi(A)}\geq2^\mu>\lambda\).
\end{hblemma}

The cardinal arithmetic is the whole hypothesis: \(\mu\leq\lambda\) because
\(2^\lambda>\lambda\), and every \(\nu<\mu\) has \(2^\nu\leq\lambda\), so the
tree \(2^{<\mu}\) has at most \(\lambda\) nodes while it has \(2^\mu>\lambda\)
branches.  This is \cite[Lemma~2.3, p.~299]{Shelah1971Stability}.

\begin{hbdefinition}{def:p03-w03-p03-032-cut}{Strongly cut sequences}
Let \(I\) be a set and \(\mathcal G\) a family of subsets of \(I\).  The
family \(\mathcal G\) \emph{strongly cuts} a sequence
\((a_n)_{n<\omega}\) of members of \(I\) when there are
\(X_k\in\mathcal G\) for \(k<\omega\) with
\[
  a_i\in X_k\quad\Longleftrightarrow\quad i<k
  \qquad(i,k<\omega).
\]
The family \(\{I\setminus X:X\in\mathcal G\}\) is written
\(I\ominus\mathcal G\).
\end{hbdefinition}

The terminology and the notation for the complementary family are those of
Erd\H{o}s and Makkai, who write \(A(-)G\) \cite[Theorem]{ErdosMakkai1966}.

\begin{hbtheorem}{thm:p03-w03-p03-032-erdos-makkai}{proved=annex:p03-w03-p03-032-erdos-makkai}{Erd\H{o}s--Makkai theorem}
Let \(m\) be an infinite cardinal, \(I\) a set with \(\abs I\leq m\), and
\(\mathcal G\) a family of subsets of \(I\) with
\(\abs{\mathcal G}>m\).  Then \(\mathcal G\) or \(I\ominus\mathcal G\)
strongly cuts some sequence of members of \(I\).
\end{hbtheorem}

For \(m=\abs I\) this is the theorem as published, for an infinite set
\(I\) and a family of more than \(\abs I\) subsets
\cite[Theorem]{ErdosMakkai1966}; the proof there, which is the one given in
the annex, uses only \(\abs I\leq m\) and \(\abs{\mathcal G}>m\).

\begin{hblemma}{thm:p03-w03-p03-032-many-to-order}{proved=annex:p03-w03-p03-032-many-to-order}{Many \(\varphi\)-types give the order property}
If \(\abs{S_\varphi(A)}>\abs A+\aleph_0\) for some small \(A\), then
\(\varphi(x;y)\) has the order property in \(T\): there are tuples
\((c_i)_{i<\omega}\) and \((b_j)_{j<\omega}\) in \(\mathfrak C\) with
\(\mathfrak C\entails\varphi(c_i;b_j)\) exactly when \(i<j\).
\end{hblemma}

This is \cite[Lemma~2.1, p.~298]{Shelah1971Stability}.

\begin{hblemma}{thm:p03-w03-p03-032-order-to-tree}{proved=annex:p03-w03-p03-032-order-to-tree}{The order property gives trees of every height}
If \(\varphi(x;y)\) has the order property in \(T\), then
\(\mathrm R_\varphi(x=x)=\infty\).
\end{hblemma}

The paper proves this through a lexicographic order on binary sequences and
compactness \cite[Lemma~2.2, pp.~298--299]{Shelah1971Stability}; the annex
places the nodes of the tree at dyadic rationals instead, which gives the
finite trees directly.

\subsection{The spectrum theorem}

\begin{hbtheorem}{thm:p03-w03-p03-032-formula}{proved=annex:p03-w03-p03-032-formula}{Equivalent forms of instability for one formula}
For a formula \(\varphi(x;y)\) of \(\Lang\) the following are equivalent.
\begin{enumerate}
  \item \(\varphi(x;y)\) has the order property in \(T\).
  \item \(\mathrm R_\varphi(x=x)=\infty\).
  \item For every infinite \(\lambda\) there is \(A\) with
        \(\abs A\leq\lambda\) and \(\abs{S_\varphi(A)}>\lambda\).
  \item There is a set \(A\) with \(\abs{S_\varphi(A)}>\abs A+\aleph_0\).
\end{enumerate}
When they fail, \(\abs{S_\varphi(A)}\leq\abs A+\aleph_0\) for every small
set \(A\).
\end{hbtheorem}

These are the properties (C), (E)/(F), (A) and (B) of the paper and their
equivalence \cite[Theorem~2.9, p.~303]{Shelah1971Stability}.

\begin{hbtheorem}{thm:p03-w03-p03-032-spectrum}{proved=annex:p03-w03-p03-032-spectrum}{Stability spectrum}
Let \(T\) be complete.  The following are equivalent.
\begin{enumerate}
  \item \(T\) is stable, that is, no formula has the order property
        (\cref{def:p03-stability}).
  \item \(T\) is \(\lambda\)-stable for some infinite cardinal \(\lambda\).
  \item \(T\) is \(\lambda\)-stable for every cardinal \(\lambda\) with
        \(\lambda^{\abs T}=\lambda\).
\end{enumerate}
Moreover, if \(T\) is stable, then
\(\abs{S_1(A)}\leq(\abs A+\aleph_0)^{\abs T}\) for every small \(A\); and if
\(T\) is unstable, then \(T\) is \(\lambda\)-stable for no infinite
\(\lambda\).
\end{hbtheorem}

The cardinals named in~(3) exist: \((2^{\abs T})^{\abs T}=2^{\abs T}\), and
more generally \((\mu^{\abs T})^{\abs T}=\mu^{\abs T}\) for every \(\mu\).
The equivalence of (1) to (3) with the order property in the form
\(\psi(a_k;a_l)\Leftrightarrow k<l\) for a single sequence is the paper's
\cite[Theorem~2.13, pp.~304--305]{Shelah1971Stability}; its condition~(2),
instability in one cardinal \(\lambda=\lambda^{\abs T}\), is the
contrapositive of (1)\(\Rightarrow\)(3) above.

\begin{hbwarning}{def:p03-w03-p03-032-arithmetic}{The hypothesis \(\lambda^{\abs T}=\lambda\) cannot be dropped}
Condition~(3) of \cref{thm:p03-w03-p03-032-spectrum} is not a claim about
every \(\lambda\geq\abs T\).  If \(\lambda=\sup_{n<\omega}\lambda_n\) with
every \(\lambda_n<\lambda\), for instance \(\lambda=\beth_\omega\), then
\(\lambda^{\aleph_0}>\lambda\) (\cref{thm:p03-w03-p03-032-koenig}), so
\(\lambda^{\abs T}>\lambda\).  The stable theory \(T_E\) of
\cref{thm:p03-w03-p03-032-te} is \(\lambda\)-stable exactly when
\(\lambda^{\aleph_0}=\lambda\); it is therefore not
\(\beth_\omega\)-stable.  Stability alone guarantees only the cardinals of
condition~(3).
\end{hbwarning}

\begin{hbwarning}{def:p03-w03-p03-032-superstable-definitions}{Two definitions of superstability}
The paper calls \(T\) \emph{superstable} when \(T\) is \(\lambda\)-stable for
every \(\lambda\geq2^{\abs T}\) \cite[\S~0, p.~273]{Shelah1971Stability}.
\Cref{def:p03-stability-spectrum} instead calls a stable theory superstable
when every complete type of a finite tuple over a set \(B\) does not fork over
a finite subset of \(B\).  The two definitions agree for every complete
theory: \cref{thm:p03-w03-p03-032-superstable} proves that each implies the
other, passing in both directions through finiteness of the degree of
\cref{def:p03-w03-p03-032-degree}.  The agreement is a theorem, not part of
either definition, and it uses stability: an unstable theory satisfies
neither condition, the first by \cref{thm:p03-w03-p03-032-spectrum} and the
second because \cref{def:p03-stability-spectrum} applies the forking
condition to stable theories only.  For the two theories of
\cref{sec:p03-w03-p03-032-separating} both conditions are also verified
directly.
\end{hbwarning}

\begin{hbexample}{def:p03-w03-p03-032-worked}{Reading the rank on equality}
Let \(T\) be the theory of an infinite set in the language with equality
only, and \(\varphi(x;y)\) the formula \(x=y\).  Above the empty set of
formulas there is a tree of height \(1\): any \(b_\varnothing\) works, since
both \(x=b_\varnothing\) and \(x\neq b_\varnothing\) are consistent.  There
is no tree of height \(2\): along the branch \(\eta=(1,1)\) the set
\(\{x=b_\varnothing,\ x=b_{(1)}\}\) forces \(b_{(1)}=b_\varnothing\), while
along \(\eta=(1,0)\) the set \(\{x=b_\varnothing,\ x\neq b_{(1)}\}\) forces
\(b_{(1)}\neq b_\varnothing\).  So \(\mathrm R_\varphi(x=x)=1\).  For a
complete \(\varphi\)-type \(p\) over \(A\) containing \(x=a\), the rank is
\(0\) and \(p_0=\{x=a\}\) decides every instance:
\(\mathrm R_\varphi(\{x=a,\ x=b\})\geq0\) exactly when \(b=a\).  For the
type \(p\) of an element outside \(A\), the rank is \(1\), \(p_0=\varnothing\)
works, and \(x=b\) belongs to \(p\) for no \(b\), in agreement with
\(\mathrm R_\varphi(\{x=b\})=0<1\).  Hence
\(\abs{S_\varphi(A)}=\abs A+1\), within the bound of
\cref{thm:p03-w03-p03-032-local-character}.
\end{hbexample}

\section{Superstability and the spectrum above \texorpdfstring{\(2^{\abs T}\)}{2 to the |T|}}
\label{sec:p03-w03-p03-032-superstable}
\index{superstability!cardinal spectrum}

The spectrum theorem leaves open what happens at cardinals \(\lambda\) with
\(\lambda^{\abs T}>\lambda\).  Shelah answers this with an ordinal-valued
degree: a formula has large degree when it can be split, again and again,
into many formulas any fixed finite number of which are jointly
inconsistent.  Shelah reads stability in every \(\lambda\geq2^{\abs T}\) off
this degree; below it carries that cardinal condition to finite local
character of forking and back.  Throughout,
\(x\) is a tuple of \(m\geq1\) variables, ``formula'' means a formula in
\(x\) with parameters from \(\mathfrak C\), and \(\kappa_T\defeq\abs T^+\).

\begin{hbdefinition}{def:p03-w03-p03-032-degree}{Degree}
A family \((\psi(x;b_l):l\in I)\) of formulas is \emph{\(k\)-inconsistent},
for a positive integer \(k\), when \(\{\psi(x;b_l):l\in w\}\) is not
realized in \(\mathfrak C\) for any \(w\subseteq I\) with \(k\) elements.
For a formula \(\theta(x)\) and an ordinal \(\alpha\), the relation
\(\operatorname{Deg}\theta\geq\alpha\) is defined by induction on
\(\alpha\).
\begin{enumerate}
  \item \(\operatorname{Deg}\theta\geq0\) iff \(\theta\) is realized in
        \(\mathfrak C\).
  \item For a limit ordinal \(\delta\), \(\operatorname{Deg}\theta\geq\delta\)
        iff \(\operatorname{Deg}\theta\geq\alpha\) for every \(\alpha<\delta\).
  \item \(\operatorname{Deg}\theta\geq\alpha+1\) iff there are a formula
        \(\psi(x;y)\) of \(\Lang\), a positive integer \(k\) and tuples
        \(b_l\) for \(l<\kappa_T\) such that
        \(\operatorname{Deg}(\theta\wedge\psi(x;b_l))\geq\alpha\) for every
        \(l\) and the family \((\psi(x;b_l):l<\kappa_T)\) is
        \(k\)-inconsistent.
\end{enumerate}
The \emph{degree} \(\operatorname{Deg}\theta\) is \(-1\) if \(\theta\) is
not realized, \(\infty\) if \(\operatorname{Deg}\theta\geq\alpha\) for every
ordinal \(\alpha\), and otherwise the ordinal \(\alpha\) with
\(\operatorname{Deg}\theta\geq\alpha\) and not
\(\operatorname{Deg}\theta\geq\alpha+1\); ordinals are below \(\infty\).
For a set \(\pi\) of formulas,
\(\operatorname{Deg}\pi\defeq\min\{\operatorname{Deg}\bigwedge\pi_0:
\pi_0\subseteq\pi\text{ finite}\}\), where \(\bigwedge\varnothing\) is
\(x=x\).  Write \(\operatorname{Deg}_m(x=x)\) when the length \(m\) of \(x\)
matters.
\index{degree of a formula}
\end{hbdefinition}

This is the paper's \(\operatorname{Deg}^m\)
\cite[Definitions~6.1--6.3, pp.~347--348]{Shelah1971Stability}, where the
\(k\)-inconsistency of clause~(3) is stated for every subset of at least
\(k\) indices, which comes to the same.  The tuple of the \(b_l\) has
\(\kappa_T\) entries, so it is small, because the saturation cardinal of
\(\mathfrak C\) exceeds \(2^{\abs T}\).  The well-definedness of the degree
is part of the next lemma.

\begin{hblemma}{thm:p03-w03-p03-032-degree-basics}{proved=annex:p03-w03-p03-032-degree-basics}{Basic properties of the degree}
Let \(\theta\), \(\theta'\), \(\theta_1\), \(\theta_2\) be formulas in \(x\).
\begin{enumerate}
  \item If \(\mathfrak C\entails\forall x\,(\theta\to\theta')\), then
        \(\operatorname{Deg}\theta\leq\operatorname{Deg}\theta'\).
  \item If \(\operatorname{Deg}\theta\geq\beta\) and \(\alpha\leq\beta\),
        then \(\operatorname{Deg}\theta\geq\alpha\).  Hence
        \(\operatorname{Deg}\theta\) is well defined.
  \item If \(\theta(x;a)\) is a formula of \(\Lang\) with parameter tuple
        \(a\) and \(a\equiv_\varnothing a'\), then
        \(\operatorname{Deg}\theta(x;a)=\operatorname{Deg}\theta(x;a')\).
  \item \(\operatorname{Deg}(\theta_1\vee\theta_2)
        =\max(\operatorname{Deg}\theta_1,\operatorname{Deg}\theta_2)\).
  \item There is an ordinal \(\alpha_0<(2^{\abs T})^+\) such that, for every
        formula \(\theta\) in \(x\), \(\operatorname{Deg}\theta\geq\alpha_0\)
        implies \(\operatorname{Deg}\theta=\infty\).
\end{enumerate}
\end{hblemma}

Items (1), (3) and~(5) are the paper's Lemma~6.1~(A), (B) and~(D)
\cite[p.~348]{Shelah1971Stability}.  Item~(4) is what makes the degree
usable for forking, whose definition is closed under finite disjunctions.

\begin{hblemma}{thm:p03-w03-p03-032-dividing-degree}{proved=annex:p03-w03-p03-032-dividing-degree}{Dividing lowers the degree}
Let \(F\) be a small set, \(\theta(x)\) a formula with parameters from \(F\)
and \(\operatorname{Deg}\theta<\infty\).  If \(\varphi(x;e)\) divides over
\(F\), then
\(\operatorname{Deg}(\theta\wedge\varphi(x;e))<\operatorname{Deg}\theta\).
\end{hblemma}

The paper records this for sequences of conjugates any \(r\) of which are
inconsistent \cite[Remark after Lemma~6.3, p.~349]{Shelah1971Stability};
dividing (\cref{def:p03-dividing-forking}) supplies such a sequence.

\begin{hblemma}{thm:p03-w03-p03-032-degree-tree}{proved=annex:p03-w03-p03-032-degree-tree}{Infinite degree gives inconsistency trees}
Suppose \(\operatorname{Deg}_m(x=x)=\infty\).  There are formulas
\(\psi_n(x;y_n)\) of \(\Lang\) and integers \(k_n\geq2\), for \(n\geq1\),
such that for every infinite small cardinal \(\lambda\) there are tuples
\(a_\eta\), for \(\eta\in\lambda^{<\omega}\setminus\{\varnothing\}\), of the
length of \(y_{\abs\eta}\), with:
\begin{enumerate}
  \item for every \(\nu\in\lambda^\omega\), the set
        \(\{\psi_n(x;a_{\nu\mathbin{\upharpoonright}n}):n\geq1\}\) is
        realized in \(\mathfrak C\);
  \item for every \(n<\omega\) and \(\eta\in\lambda^n\), the family
        \((\psi_{n+1}(x;a_{\eta\frown\langle i\rangle}):i<\lambda)\) is
        \(k_{n+1}\)-inconsistent.
\end{enumerate}
\end{hblemma}

This is the passage from (A) to (B) of
\cite[Theorem~6.4, pp.~349--351]{Shelah1971Stability} followed by the
compactness step of \cite[proof of Theorem~6.5,
pp.~353--354]{Shelah1971Stability}.  The annex replaces the paper's
induction on the levels by a pruning of finite trees, which yields the same
formulas \(\psi_n\) at every node of a level.

\begin{hblemma}{thm:p03-w03-p03-032-dividing-tree-count}{proved=annex:p03-w03-p03-032-dividing-tree-count}{Inconsistency trees give many types}
Let \(\lambda\) be a small cardinal with
\(\lambda^{\aleph_0}>\lambda\geq2^{\aleph_0}\), and let \(\psi_n\),
\(k_n\) and \((a_\eta)\) be as in
\cref{thm:p03-w03-p03-032-degree-tree} for this \(\lambda\).  Let \(A\) be
the set of entries of the tuples \(a_\eta\).  Then \(\abs A\leq\lambda\)
and \(\abs{S_m(A)}>\lambda\); consequently \(T\) is not
\(\lambda\)-stable.
\end{hblemma}

The paper proves the count for trees in which any two siblings are
inconsistent, where it holds for every \(\lambda\) with
\(\lambda^{\aleph_0}>\lambda\) \cite[Theorem~6.5,
pp.~353--354]{Shelah1971Stability}.  With \(k_n\)-inconsistent siblings a
type can contain up to \(2^{\aleph_0}\) branches, and the hypothesis
\(\lambda\geq2^{\aleph_0}\) absorbs that factor.

The converse direction, from finite degree back to the cardinal condition,
needs a count of the complete types that keep the degree of a formula.  The
count rests on a supply of long sequences on which finitely many formulas
behave uniformly, and that supply rests on a pressing-down argument.

\begin{hblemma}{thm:p03-w03-p03-032-pressing-down}{proved=annex:p03-w03-p03-032-pressing-down}{A regressive function is bounded on a large set}
Let \(\lambda\) be a regular uncountable cardinal and \(f\) a function
defined on the nonzero limit ordinals \(k<\lambda\) with \(f(k)<k\) for each of
them.  Then there is \(j<\lambda\) such that \(f(k)\leq j\) holds for
\(\lambda\) many nonzero limit ordinals \(k<\lambda\).
\end{hblemma}

This weak form of the pressing-down lemma is all that is used below; the
stationary-set form is \cref{thm:p06-fodor}.

\begin{hblemma}{thm:p03-w03-p03-032-uniform}{proved=annex:p03-w03-p03-032-uniform}{Uniform subsequences in a stable theory}
Let \(T\) be stable, \(\lambda\) a regular uncountable small cardinal,
\(D\) a set with \(\abs D<\lambda\), \(N\geq1\), and \(\Delta\) a finite
set of formulas \(\delta(y_0,\ldots,y_{N-1})\) with parameters from \(D\),
where \(y_0,\ldots,y_{N-1}\) are tuples of variables of one length \(l\).
For every sequence \((b_k:k<\lambda)\) of \(l\)-tuples there is
\(K\subseteq\lambda\) with \(\abs K=\lambda\) such that, for each
\(\delta\in\Delta\), either
\(\mathfrak C\entails\delta(b_{k_0},\ldots,b_{k_{N-1}})\) for all
\(k_0<\cdots<k_{N-1}\) in \(K\), or for none of them.
\end{hblemma}

This is the sequence form of the paper's theorem on indiscernible sets
\cite[Theorem~5.8\,(A), p.~340]{Shelah1971Stability}, which asks in addition
that the truth value not depend on the order of the indices; the paper does
not prove that theorem but refers to earlier work for it
\cite[Remark after Theorem~5.8, p.~340]{Shelah1971Stability}.  The annex
proves the sequence form from the finite local character of stable formulas
(\cref{thm:p03-w03-p03-032-local-character}); stability enters only there.
The lemma gives \(\lambda\) indices out of \(\lambda\), with \(\lambda\)
regular and uncountable; for finite colourings of \(N\)-tuples of a
countable sequence, an infinite uniform subsequence exists in every theory by
Ramsey's theorem, which is how
\cref{thm:p03-w03-p03-035-em-extraction} is obtained.

\begin{hblemma}{thm:p03-w03-p03-032-same-degree}{proved=annex:p03-w03-p03-032-same-degree}{Few complete extensions of the same degree}
Let \(T\) be stable, \(A\) a small set, \(x\) a tuple of \(m\) variables,
\(\theta(x)\) a formula with parameters from \(A\) and
\(\alpha=\operatorname{Deg}\theta<\infty\).  Put
\[
  P_\theta\defeq\{p\in S_m(A):\theta\in p,\ \operatorname{Deg}p=\alpha\}.
\]
\begin{enumerate}
  \item For every formula \(\varphi(x;y)\) of \(\Lang\), the set
        \(\{p\mathbin{\upharpoonright}\varphi:p\in P_\theta\}\) has at most
        \(\abs T\) members.
  \item \(\abs{P_\theta}\leq2^{\abs T}\).
\end{enumerate}
\end{hblemma}

Item~(1) is the paper's Theorem~6.7, which excludes \(\abs T^+\) types of
the same degree that are pairwise contradictory on one formula
\cite[Theorem~6.7, pp.~355--356]{Shelah1971Stability}; item~(2) is the
bound of its Corollary~6.8\,(B) \cite[p.~357]{Shelah1971Stability}.  The
annex follows the plan of the paper's proof of Theorem~6.7: a finite part
of each local type decides it, these finite parts are put on a uniform
sequence, and a fixed alternation pattern then gives a \(k\)-inconsistent
family of \(\kappa_T\) formulas of the same degree.  Where the paper negates
a fixed block of instances, the annex pairs consecutive members of the
sequence.

\begin{hblemma}{thm:p03-w03-p03-032-forking-chain}{proved=annex:p03-w03-p03-032-forking-chain}{Infinite degree gives a type forking over every finite set}
If \(\operatorname{Deg}_m(x=x)=\infty\), there are a countable set \(B\)
and a complete type \(p\) in \(m\) variables over \(B\) that forks over
every finite subset of \(B\).
\end{hblemma}

The paper has no forking relation and no counterpart of this lemma; its
nearest statement is the tree property~(E) of
\cite[Corollary~6.10, p.~357]{Shelah1971Stability}.  The annex turns the
inconsistency trees of \cref{thm:p03-w03-p03-032-degree-tree} into a chain
of formulas each of which divides over the parameters of the earlier ones,
using the extraction of indiscernible sequences
(\cref{thm:p03-w03-p03-035-em-extraction}).

\begin{hbtheorem}{thm:p03-w03-p03-032-superstable}{proved=annex:p03-w03-p03-032-superstable}{Superstability above \(2^{\abs T}\)}
For a complete theory \(T\) the following are equivalent.
\begin{enumerate}
  \item \(T\) is \(\lambda\)-stable for every cardinal
        \(\lambda\geq2^{\abs T}\);
  \item \(T\) is \(\lambda\)-stable for some cardinal \(\lambda\) with
        \(\lambda^{\aleph_0}>\lambda\geq2^{\aleph_0}\);
  \item \(T\) is stable and \(\operatorname{Deg}_m(x=x)<\infty\) for every
        \(m\geq1\);
  \item \(T\) is superstable in the sense of
        \cref{def:p03-stability-spectrum};
  \item \(T\) is stable and \(\operatorname{Deg}_1(x=x)<\infty\).
\end{enumerate}
Under these conditions \(\abs{S_1(A)}\leq\abs A+2^{\abs T}\) for every small
set \(A\).  More precisely, under~(3), if \(p\) is a complete type in
finitely many variables over a set \(B\) and \(q\subseteq p\) is finite with
\(\operatorname{Deg}\bigwedge q=\operatorname{Deg}p\), then \(p\) does not
fork over the finite set of parameters of \(q\).
\end{hbtheorem}

In the paper, (1) is the definition of superstability
\cite[\S~0, p.~273]{Shelah1971Stability}, and for stable \(T\) it is shown
equivalent to finiteness of \(\operatorname{Deg}_1(x=x)\) and to stability
in one \(\lambda\) with \(\lambda^{\aleph_0}>\lambda\geq2^{\abs T}\)
\cite[Corollary~6.10, pp.~357--358]{Shelah1971Stability}.  The step from
(5) to (1), with the bound \(\abs A+2^{\abs T}\), is the paper's
\cite[Theorem~6.9, p.~357]{Shelah1971Stability}, whose proof chooses for
each type a finite part of the same degree, as
\cite[Lemma~6.2\,(B), p.~348]{Shelah1971Stability} permits, and counts
with \cref{thm:p03-w03-p03-032-same-degree}.  The steps from (3) to (4)
and from (4) to (3) connect the cardinal notion to the forking definition
used in this handbook; together they give the agreement recorded in
\cref{def:p03-w03-p03-032-superstable-definitions}.

The hypotheses are sharp in two places.  The exponent \(\aleph_0\) in (2)
cannot be replaced by \(\abs T\): every stable theory is \(\lambda\)-stable
for \(\lambda=\lambda^{\abs T}\) (\cref{thm:p03-w03-p03-032-spectrum}),
including the non-superstable \(T_E\) of \cref{thm:p03-w03-p03-032-te}.  And
(1) concerns only \(\lambda\geq2^{\abs T}\): the superstable theory \(T_P\)
of \cref{thm:p03-w03-p03-032-tp} is not \(\aleph_0\)-stable.

\section{Theories separating the stability classes}
\label{sec:p03-w03-p03-032-separating}

Each separating theory is analysed through one device: in it, the type of a
tuple is decided by its quantifier-free type, which reduces type counting to
counting patterns of atomic formulas.

\begin{hblemma}{thm:p03-w03-p03-032-qf-test}{proved=annex:p03-w03-p03-032-qf-test}{Quantifier-free types decide types}
Let \(T_0\) be a consistent theory in a relational language without
constant symbols.  Here a model is \emph{\(\omega\)-saturated} when every
set of formulas in one variable with parameters from a finite subset of it,
each finite subset of which is satisfied in the model, is realized in it.
For \(\omega\)-saturated models \(\mathcal M,\mathcal N\)
of \(T_0\), call a map \(f\) from a finite subset of \(M\) into \(N\) a
\emph{finite partial isomorphism} when, for every atomic formula
\(\alpha(\bar x)\) and every tuple \(\bar a\) from its domain,
\(\mathcal M\entails\alpha(\bar a)\) iff
\(\mathcal N\entails\alpha(f\bar a)\).  Suppose that for all such
\(\mathcal M,\mathcal N\), every finite partial isomorphism \(f\) and every
\(c\in M\) and \(d\in N\), some finite partial isomorphism extends \(f\) and
has \(c\) in its domain, and some has \(d\) in its range.  Then:
\begin{enumerate}
  \item finite partial isomorphisms between \(\omega\)-saturated models of
        \(T_0\) preserve every formula;
  \item \(T_0\) is complete;
  \item in a model of \(T_0\), tuples satisfying the same atomic formulas
        have the same complete type.  In particular, for tuples \(c,c'\) and
        a set \(A\) in \(\mathfrak C\), if \((c,\bar a)\) and
        \((c',\bar a)\) satisfy the same atomic formulas for every finite
        tuple \(\bar a\) from \(A\), then \(c\equiv_A c'\).
\end{enumerate}
\end{hblemma}

\subsection{Independent unary predicates: superstable, not
\texorpdfstring{\(\omega\)}{omega}-stable}

\begin{hbdefinition}{def:p03-w03-p03-032-tp}{The theory \(T_P\)}
The language has unary relation symbols \(P_n\) for \(n<\omega\).  For a
finite set \(I\subseteq\omega\) and \(\sigma:I\to\{0,1\}\), write
\(P^\sigma(x)\) for \(\bigwedge_{n\in I}P_n^{\sigma(n)}(x)\), where
\(P_n^1=P_n\) and \(P_n^0=\neg P_n\).  The theory \(T_P\) has, for every such
\(\sigma\) and every \(k\geq1\), the axiom saying that at least \(k\)
elements satisfy \(P^\sigma(x)\).  For an element \(c\) of a model, its
\emph{pattern} is \(\sigma_c\in2^\omega\) with \(\sigma_c(n)=1\) iff
\(P_n(c)\).
\end{hbdefinition}

\begin{hbtheorem}{thm:p03-w03-p03-032-tp}{proved=annex:p03-w03-p03-032-tp}{\(T_P\) is superstable and not \(\omega\)-stable}
\begin{enumerate}
  \item \(T_P\) has a model and is complete, and in \(\mathfrak C\entails T_P\)
        the type of an element over a set \(A\) is decided by whether it lies
        in \(A\) and, if not, by its pattern.
  \item For every small \(A\),
        \(\abs{S_1(A)}=\abs A+2^{\aleph_0}\).  Hence \(T_P\) is
        \(\lambda\)-stable exactly when \(\lambda\geq2^{\aleph_0}\); it is
        stable, it is not \(\omega\)-stable, and it is \(\lambda\)-stable for
        every \(\lambda\geq2^{\abs{T_P}}\).
  \item For a tuple \(a\) and small sets \(C\subseteq B\), if every entry of
        \(a\) that lies in \(B\) lies in \(C\), then
        \(\operatorname{tp}(a/B)\) does not fork over \(C\).  In particular
        \(\operatorname{tp}(a/B)\) does not fork over the finite set of
        entries of \(a\) lying in \(B\), so \(T_P\) is superstable in the
        sense of \cref{def:p03-stability-spectrum}.
\end{enumerate}
\end{hbtheorem}

\subsection{Refining equivalence relations: stable, not superstable}

\begin{hbdefinition}{def:p03-w03-p03-032-te}{The theory \(T_E\) and the height function}
The language has binary relation symbols \(E_n\) for \(n<\omega\).  The
theory \(T_E\) says: each \(E_n\) is an equivalence relation;
\(E_{n+1}\subseteq E_n\); \(E_0\) has infinitely many classes; and every
\(E_n\)-class contains infinitely many \(E_{n+1}\)-classes (one axiom for
each \(n\) and each finite lower bound).  In a model, the \emph{height} of
two elements is
\[
  h(x,y)\defeq\min\{n<\omega:\neg\,x\mathrel{E_n}y\},
  \qquad h(x,y)\defeq\omega\text{ if no such }n\text{ exists},
\]
so that \(x\mathrel{E_n}y\) holds exactly when \(n<h(x,y)\).
\end{hbdefinition}

For an infinite cardinal \(\lambda\), the structure \(\mathcal M_\lambda\)
with universe the set \(\lambda^\omega\) of \(\omega\)-sequences of ordinals
below \(\lambda\), and \(\eta\mathrel{E_n}\eta'\) iff
\(\eta\mathbin{\upharpoonright}(n+1)=\eta'\mathbin{\upharpoonright}(n+1)\), is a
model of \(T_E\); in it \(h(\eta,\eta')\) is the first place where
\(\eta\) and \(\eta'\) differ.

\begin{hblemma}{thm:p03-w03-p03-032-koenig}{proved=annex:p03-w03-p03-032-koenig}{Countable cofinality raises the power}
If \(\lambda=\sup_{n<\omega}\lambda_n\) with \(\lambda_n<\lambda\) for every
\(n\), then \(\lambda^{\aleph_0}>\lambda\).
\end{hblemma}

\begin{hbtheorem}{thm:p03-w03-p03-032-te}{proved=annex:p03-w03-p03-032-te}{\(T_E\) is stable and not superstable}
\begin{enumerate}
  \item \(h(x,z)\geq\min(h(x,y),h(y,z))\), with equality when
        \(h(x,y)\neq h(y,z)\).  \(T_E\) is complete, and in
        \(\mathfrak C\entails T_E\) the type of a tuple over a set is decided by
        the equalities among its entries and the parameters and by the
        heights between them.
  \item For every infinite \(\lambda\) and every \(A\) with
        \(\abs A\leq\lambda\),
        \(\abs{S_1(A)}\leq\lambda^{\aleph_0}\); and for every infinite
        \(\lambda\) some \(A\) with \(\abs A=\lambda\) has
        \(\abs{S_1(A)}\geq\lambda^{\aleph_0}\).  Hence \(T_E\) is
        \(\lambda\)-stable exactly when \(\lambda^{\aleph_0}=\lambda\); it is
        stable, and it is not \(\beth_\omega\)-stable although
        \(\beth_\omega\geq2^{\abs{T_E}}\).
  \item There are elements \(a_n\) \((n<\omega)\) and \(c\) with
        \(h(a_n,a_m)=\min(n,m)\) for \(n\neq m\) and \(h(c,a_n)=n\); for
        \(B=\{a_n:n<\omega\}\), the type \(\operatorname{tp}(c/B)\) forks over
        every finite subset of \(B\).  So \(T_E\) is not superstable in the
        sense of \cref{def:p03-stability-spectrum}.
\end{enumerate}
\end{hbtheorem}

\subsection{Stable and simple}

The random graph \(T_{\rm RG}\) is unstable: the edge relation has the order
property (\cref{def:p03-w03-p03-031-random-graph-example}), so by
\cref{thm:p03-w03-p03-032-spectrum} it is \(\lambda\)-stable for no infinite
\(\lambda\).  It is simple, with finite bases for local character
(\cref{thm:p03-w03-p03-031-random-graph}).  Together with
\cref{thm:p03-w03-p03-032-tp,thm:p03-w03-p03-032-te} this gives one theory
in each gap: \(T_P\) between \(\omega\)-stable and superstable, \(T_E\)
between superstable and stable, and \(T_{\rm RG}\) between stable and simple.

\begin{hbwarning}{def:p03-w03-p03-032-limits}{What the separating theories do and do not show}
Each theory witnesses that one inclusion between classes is strict; none of
them proves an inclusion.  The inclusions themselves, \(\omega\)-stable
inside superstable inside stable inside simple, are the statements cited in
\cref{thm:p03-nonforking-calculus} and \cref{sec:p03-stability}.  The
counts in \cref{thm:p03-w03-p03-032-tp,thm:p03-w03-p03-032-te} use the
countable languages of these theories: for \(T_P\) the value \(2^{\aleph_0}\)
is the number of patterns, and for \(T_E\) the exponent \(\aleph_0\) is the
number of relations \(E_n\) along a descending chain of classes.
\end{hbwarning}

\paragraph{Relationships.}
The definitions used are those of \cref{sec:p03-stability}.  The ambient
model, smallness, and the realization of small finitely satisfiable sets of
formulas are those of \cref{def:p03-w03-p03-031-monster} and
\cref{thm:p03-w03-p03-031-saturation}.  The forking laws that hold in stable
and simple theories are developed in \cref{sec:p03-w03-p03-031-forking-calculus};
the separating theories here are the examples against which those laws can be
read, and the random-graph computation there is the third separation.  Morley
rank and \(\omega\)-stability are treated in \cref{sec:p03-morley-rank}.

\paragraph{Sources.}
Shelah \cite[\S~2, pp.~296--305]{Shelah1971Stability} for the rank, the
formula-level equivalences and the spectrum theorem, and \cite[\S~0,
p.~273]{Shelah1971Stability} for the cardinal definition of superstability,
\cite[\S~6, pp.~347--358]{Shelah1971Stability} for the degree and its
relation to superstability, with \cite[Theorem~5.8, p.~340]{Shelah1971Stability}
for the uniform sequences behind the count;
Erd\H{o}s and Makkai \cite[Theorem]{ErdosMakkai1966} for the combinatorial
theorem.  The theories \(T_P\) and \(T_E\), their analysis, the dyadic
tree of \cref{thm:p03-w03-p03-032-order-to-tree}, the proofs of
\cref{thm:p03-w03-p03-032-pressing-down,thm:p03-w03-p03-032-uniform}, and the
passage from infinite degree to forking in
\cref{thm:p03-w03-p03-032-forking-chain} are presented here with full
proofs.
